% Author: Grayson Orr
% Course: IN721: Mobile Application Development

\documentclass{article}
\author{}

\usepackage{graphicx}
\usepackage{wrapfig}
\usepackage{enumerate}
\usepackage{hyperref}
\usepackage[margin = 2.25cm]{geometry}
\usepackage[table]{xcolor}
\usepackage{fancyhdr}
\hypersetup{
  colorlinks = true,
  urlcolor = blue
}
\setlength\parindent{0pt}
\pagestyle{fancy}
\fancyhf{}
\rhead{College of Engineering, Construction and Living Sciences\\Bachelor of Information Technology}
\lfoot{Presentation\\Version 1, Semester Two, 2021}
\rfoot{\thepage}
 
\begin{document}

\begin{figure}
	\centering
	\includegraphics[width=50mm]{./resources/img/logo.png}
\end{figure}

\title{College of Engineering, Construction and Living Sciences\\Bachelor of Information Technology\\IN721: Mobile Application Development\\Level 7, Credits 15\\\textbf{Presentation}}
\date{}
\maketitle

\section*{Assessment Overview}
In this assessment, you will present a mobile-related topic to the course lecturer. The information presented must be in a \textbf{README.md} file. Also, you need to provide a code example to accompany the README.md file. The main purpose of this assessment is to demonstrate your ability to identify and effectively articulate an intermediate/advanced topic in \textbf{Android}. 

\section*{Learning Outcomes}
At the successful completion of this course, learners will be able to:
\begin{enumerate}
	\item Implement \& publish complete, non-trivial, industry-standard mobile applications following sound architectural \& code-quality standards.
	\item Identify relevant use cases for a mobile computing scenario \& incorporate them into an effective user experience design.
	\item Follow industry standard software engineering practice in the design of mobile applications.
\end{enumerate}

\section*{Assessment Table}
\renewcommand{\arraystretch}{1.5}
\begin{tabular}{|l|l|l|l|l|}
	\hline
	\vtop{\hbox{\strut \textbf{Assessment}}\hbox{\strut \textbf{Activity}}} & \textbf{Weighting} & \vtop{\hbox{\strut \textbf{Learning}}\hbox{\strut \textbf{Outcomes}}} & \vtop{\hbox{\strut \textbf{Assessment}}\hbox{\strut \textbf{Grading Scheme}}} & \vtop{\hbox{\strut \textbf{Completion}}\hbox{\strut \textbf{Requirements}}} \\
	
	\hline
	
	\small Practicals                                                        & \small 10\%        & \small 2, 3                                                           & \small CRA                                                                    & \small Cumulative                                                           \\ \hline
	\small Project                                                          & \small 70\%        & \small 1, 2, 3                                                        & \small CRA                                                                    & \small Cumulative                                                           \\ \hline
	\small Presentation                                                          & \small 20\%        & \small 2, 3                                                        & \small CRA                                                                    & \small Cumulative                                                           \\ \hline
\end{tabular}

\section*{Conditions of Assessment}
You will complete this assessment during your learner managed time, however, there will be availability during the teaching sessions to discuss the requirements \& your progress of this assessment. This assessment will need to be completed by \textbf{Thursday, 18 November 2021} at \textbf{5 PM}.

\section*{Pass Criteria}
This assessment is criterion-referenced (CRA) with a cumulative pass mark of \textbf{50\%} over all assessments in \textbf{IN721: Mobile Application Development}.

\section*{Authenticity}
All parts of your submitted assessment \textbf{must} be completely your work \& any references \textbf{must} be cited appropriately including, externally-sourced graphic elements. Provide your references in a \textbf{README.md} file. All media \textbf{must} be royalty free (or legally purchased) for educational use. Failure to do this will result in a mark of \textbf{zero} for this assessment.

\section*{Policy on Submissions, Extensions, Resubmissions \& Resits}
The school's process concerning submissions, extensions, resubmissions \& resits complies with \textbf{Otago Polytechnic} policies. Learners can view policies on the \textbf{Otago Polytechnic} website located at \href{https://www.op.ac.nz/about-us/governance-and-management/policies}{https://www.op.ac.nz/about-us/governance-and-management/policies}.

\section*{Submissions}
You \textbf{must} submit all presentation files via \textbf{GitHub Classroom}. Here is the URL to the repository you will use for your submission – \href{https://classroom.github.com/a/EQ7POFAY}{https://classroom.github.com/a/EQ7POFAY}. The latest presentation files in the \textbf{main} branch will be used to mark against the \textbf{Documentation} criterion. Partial marks \textbf{are not} given for functionality in other branches. Late submissions will incur a \textbf{10\% penalty per day}, rolling over at \textbf{5:00 PM}.

\section*{Extensions}
Familiarise yourself with the assessment due date. If you need an extension, contact the course lecturer before the due date. If you require more than a week's extension, a medical certificate or support letter from your manager may be needed.

\section*{Resubmissions}
Learners may be requested to resubmit an assessment following a rework of part/s of the original assessment. Resubmissions are to be completed within a negotiable short time frame \& usually \textbf{must} be completed within the timing of the course to which the assessment relates. Resubmissions will be available to learners who have made a genuine attempt at the first assessment opportunity \& achieved a \textbf{D grade (40-49\%)}. The maximum grade awarded for resubmission will be \textbf{C-}.

\section*{Resits}
Resits \& reassessments \textbf{are not} applicable in \textbf{IN721: Mobile Application Development}. 

\section*{Instructions}

A list of topics are available in the \textbf{Microsoft Teams} course channel, under the \textbf{Files} tab. The file is called \textbf{presentation-topics.docx}.

\subsection*{Documentation - Learning Outcomes 2, 3 (50\%)}
\begin{itemize}
	\item Documentation must contain the following sections:
	\begin{itemize}
		\item Overview - a brief description of what the topic is.
		\item Dependencies - it may include the operating system, name, version number, etc.
		\item Code example - a description of each code snippet in relation to the topic. It means you \textbf{only} have to describe the essential files.
		\item References - the information in your documentation is referenced using \textbf{APA 7th edition}.
		\begin{itemize}
			\item \textbf{Resource:} \href{https://studentservices.op.ac.nz/learning-support/citingandreferencing}{https://studentservices.op.ac.nz/learning-support/citingandreferencing}
		\end{itemize}
	\end{itemize}
	\item Use of \textbf{Markdown}, i.e., bold text, code blocks, etc.
\end{itemize}

\subsection*{Presentation 2, 3 (50\%)}
\begin{itemize}
	\item Present your documentation, i.e., \textbf{README.md} to the course lecturer, or record your presentation. If you wish to record your presentation, you must:
	\begin{itemize}
		\item Upload your presentation to your OP student OneDrive.
		\item Provide a link to your presentation in your documentation.
	\end{itemize}
	\item Answer the following:
	\begin{itemize}
		\item Describe how would you implement it into your travelling \textbf{Project}?
	\end{itemize}
\end{itemize}
\end{document}
